\documentclass[10pt,a4paper]{amsart}
\usepackage{amsmath,amssymb,amsthm,graphicx}
\usepackage{verbatim}
\newenvironment{reference}{\comment}{\endcomment}
%\newenvironment{reference}{}{}
\newenvironment{slogan}{\comment}{\endcomment}
\newenvironment{history}{\comment}{\endcomment}

% For commutative diagrams we use Xy-pic
\usepackage[all]{xy}

% We use 2cell for 2-commutative diagrams.
\xyoption{2cell}
% Theorem environments.
%
\theoremstyle{plain}
\newtheorem{theorem}[subsection]{Theorem}
\newtheorem{proposition}[subsection]{Proposition}
\newtheorem{lemma}[subsection]{Lemma}

\theoremstyle{definition}
\newtheorem{definition}[subsection]{Definition}
\newtheorem{example}[subsection]{Example}
\newtheorem{exercise}[subsection]{Exercise}
\newtheorem{situation}[subsection]{Situation}

\theoremstyle{remark}
\newtheorem{remark}[subsection]{Remark}
\newtheorem{remarks}[subsection]{Remarks}

\numberwithin{equation}{subsection}

% Macros
%
\def\lim{\mathop{\mathrm{lim}}\nolimits}
\def\colim{\mathop{\mathrm{colim}}\nolimits}
\def\Spec{\mathop{\mathrm{Spec}}}
\def\Hom{\mathop{\mathrm{Hom}}\nolimits}
\def\Ext{\mathop{\mathrm{Ext}}\nolimits}
\def\SheafHom{\mathop{\mathcal{H}\!\mathit{om}}\nolimits}
\def\SheafExt{\mathop{\mathcal{E}\!\mathit{xt}}\nolimits}
\def\Sch{\mathit{Sch}}
\def\Mor{\mathop{\mathrm{Mor}}\nolimits}
\def\Ob{\mathop{\mathrm{Ob}}\nolimits}
\def\Sh{\mathop{\mathit{Sh}}\nolimits}
\def\NL{\mathop{N\!L}\nolimits}
\def\CH{\mathop{\mathrm{CH}}\nolimits}
\def\proetale{{pro\text{-}\acute{e}tale}}
\def\etale{{\acute{e}tale}}
\def\QCoh{\mathit{QCoh}}
\def\Ker{\mathop{\mathrm{Ker}}}
\def\Im{\mathop{\mathrm{Im}}}
\def\Coker{\mathop{\mathrm{Coker}}}
\def\Coim{\mathop{\mathrm{Coim}}}

% Boxtimes
%
\DeclareMathSymbol{\boxtimes}{\mathbin}{AMSa}{"02}

%
% Macros for moduli stacks/spaces
%
\def\QCohstack{\mathcal{QC}\!\mathit{oh}}
\def\Cohstack{\mathcal{C}\!\mathit{oh}}
\def\Spacesstack{\mathcal{S}\!\mathit{paces}}
\def\Quotfunctor{\mathrm{Quot}}
\def\Hilbfunctor{\mathrm{Hilb}}
\def\Curvesstack{\mathcal{C}\!\mathit{urves}}
\def\Polarizedstack{\mathcal{P}\!\mathit{olarized}}
\def\Complexesstack{\mathcal{C}\!\mathit{omplexes}}
% \Pic is the operator that assigns to X its picard group, usage \Pic(X)
% \Picardstack_{X/B} denotes the Picard stack of X over B
% \Picardfunctor_{X/B} denotes the Picard functor of X over B
\def\Pic{\mathop{\mathrm{Pic}}\nolimits}
\def\Picardstack{\mathcal{P}\!\mathit{ic}}
\def\Picardfunctor{\mathrm{Pic}}
\def\Deformationcategory{\mathcal{D}\!\mathit{ef}}
\usepackage{cleveref}
\crefname{originalreference}{original reference}{original references}
\Crefname{originalreference}{Original reference}{Original references}
\makeatletter
\newlabel{situation-setup}{{02QL}{source}{}{sourcepointer-1}{}}
\newlabel{situation-setup@cref}{{[originalreference][0][]02QL}{[1][0][]source}{}{}{}}
\makeatother
\crefname{originalsection}{original section}{original sections}
\Crefname{originalsection}{Original section}{Original sections}
\makeatletter
\newlabel{section-intersection-regular}{{0FEX}{source}{}{sourcepointer-2}{}}
\newlabel{section-intersection-regular@cref}{{[originalsection][0][]0FEX}{[1][0][]source}{}{}{}}
\makeatother
\crefname{originallemma}{original lemma}{original lemmas}
\Crefname{originallemma}{Original lemma}{Original lemmas}
\makeatletter
\newlabel{more-morphisms-lemma-morphism-regular-schemes-is-lci}{{0E9K}{source}{}{sourcepointer-3}{}}
\newlabel{more-morphisms-lemma-morphism-regular-schemes-is-lci@cref}{{[originallemma][0][]0E9K}{[1][0][]source}{}{}{}}
\makeatother
\crefname{originallemma}{original lemma}{original lemmas}
\Crefname{originallemma}{Original lemma}{Original lemmas}
\makeatletter
\newlabel{lemma-lci-gysin-easy}{{0FF7}{source}{}{sourcepointer-4}{}}
\newlabel{lemma-lci-gysin-easy@cref}{{[originallemma][0][]0FF7}{[1][0][]source}{}{}{}}
\makeatother
\crefname{originallemma}{original lemma}{original lemmas}
\Crefname{originallemma}{Original lemma}{Original lemmas}
\makeatletter
\newlabel{lemma-lci-gysin-commutes}{{0FF6}{source}{}{sourcepointer-5}{}}
\newlabel{lemma-lci-gysin-commutes@cref}{{[originallemma][0][]0FF6}{[1][0][]source}{}{}{}}
\makeatother
\title{The Gysin map of a morphism between quasi-compact regular schemes with affine diagonal is multiplicative in the regular intersection theory}
\author{The Stacks Project Authors}
\date{stacks-project:chow.tex:local-snapshot-20261012; GNU FDL 1.2 or later}
\begin{document}
\maketitle
\noindent\textit{Editorial note.} All delivered mathematical text is the author's own, copied from the pinned Stacks Project chapter file chow.tex; no statement or proof step was rewritten, reordered or removed.\par\medskip
\noindent\textit{Editorial note.} Preamble provenance: the pinned chapter file opens with an input of the shared preamble and carries no preamble text of its own, so the project's single Stacks-wide preamble staged at sources/stacks-chow/preamble.tex is excerpted verbatim as two spans, lines 14-18 (Xy-pic setup) and lines 154-217 (the native theorem declarations and mathematical macros). Nothing was added, removed or redefined; the wrapper class and the amsmath/amssymb/amsthm/graphicx base packages are supplied by the assembler.\par\medskip
\noindent\textit{Editorial note.} The counted statement and proof contain no cross reference, no citation and no included file, so the delivered body is self contained.\par\medskip
\noindent\textit{Editorial note.} Canonical location of the counted proposition: Stacks tag 0FF9 (sources/stacks-chow/pinned-original-source.tex line 13868).\par\medskip
\noindent\textit{Allowed standard prerequisites (editorial).}\begin{itemize}
\item A morphism between quasi-compact regular schemes with affine diagonal is local complete intersection. More on Morphisms, Lemma labelled more-morphisms-lemma-morphism-regular-schemes-is-lci, cited at proof lines 13882-13884. Gives the first assertion and the availability of f\textasciicircum{}!.
\item The value of the Gysin map on the fundamental class. Same chapter, Lemma labelled lemma-lci-gysin-easy, cited at proof lines 13886-13888. Supplies f\textasciicircum{}![Y] = [X].
\item The intersection product on regular schemes via the Chern character. Same chapter, Section labelled section-intersection-regular, invoked at proof lines 13888-13896 and 13904-13910. Writes the Chow classes as Chern characters against fundamental classes.
\item Commutation of the Gysin map with bivariant classes. Same chapter, Lemma labelled lemma-lci-gysin-commutes, cited at proof lines 13898-13902. Supplies the middle equalities of the chain.
\end{itemize}\medskip
\ifdefined\hypertarget\hypertarget{sourcepointer-1}{}\fi
\noindent\textit{Source reference.} Original reference 02QL: sources/stacks-chow/pinned-original-source.tex, source lines 1372--1372. Same-chapter pointer into the author environment situation at source line 1372; invoked at proof line None of the counted proof. It is retained only as a navigation pointer; every load-bearing step is inside the delivered span.\par\medskip
\ifdefined\hypertarget\hypertarget{sourcepointer-2}{}\fi
\noindent\textit{Source reference.} Original section 0FEX: sources/stacks-chow/pinned-original-source.tex, source lines 13290--13290. Same-chapter pointer into the author environment lemma at source line 13290; invoked at proof line 13889 of the counted proof. It is retained only as a navigation pointer; every load-bearing step is inside the delivered span.\par\medskip
\ifdefined\hypertarget\hypertarget{sourcepointer-3}{}\fi
\noindent\textit{Source reference.} Original lemma 0E9K: sources/stacks-chow/tags.txt, source lines 17855--17855. Cross-chapter standard prerequisite located in the project tag list staged in the local snapshot; invoked at proof line 13885 of the counted proof. It is retained only as a navigation pointer; every load-bearing step is inside the delivered span.\par\medskip
\ifdefined\hypertarget\hypertarget{sourcepointer-4}{}\fi
\noindent\textit{Source reference.} Original lemma 0FF7: sources/stacks-chow/pinned-original-source.tex, source lines 13643--13643. Same-chapter pointer into the author environment lemma at source line 13643; invoked at proof line 13886 of the counted proof. It is retained only as a navigation pointer; every load-bearing step is inside the delivered span.\par\medskip
\ifdefined\hypertarget\hypertarget{sourcepointer-5}{}\fi
\noindent\textit{Source reference.} Original lemma 0FF6: sources/stacks-chow/pinned-original-source.tex, source lines 13613--13613. Same-chapter pointer into the author environment lemma at source line 13613; invoked at proof line 13892 of the counted proof. It is retained only as a navigation pointer; every load-bearing step is inside the delivered span.\par\medskip
\begin{lemma}
\label{lemma-lci-gysin-product-regular}
Let $(S, \delta)$ be as in Situation \ref{situation-setup}.
Let $f : X \to Y$ be a morphism of schemes locally of finite
type over $S$ such that both $X$ and $Y$ are quasi-compact,
regular, have affine diagonal, and finite dimension.
Then $f$ is a local complete intersection morphism.
Assume moreover the Gysin map exists for $f$. Then
$$
f^!(\alpha \cdot \beta) = f^!\alpha \cdot f^!\beta
$$
in $\CH^*(X) \otimes \mathbf{Q}$ where the intersection product
is as in Section \ref{section-intersection-regular}.
\end{lemma}

\begin{proof}
The first statement follows from
More on Morphisms, Lemma
\ref{more-morphisms-lemma-morphism-regular-schemes-is-lci}.
Observe that $f^![Y] = [X]$, see Lemma \ref{lemma-lci-gysin-easy}.
Write $\alpha = ch(\alpha') \cap [Y]$ and $\beta = ch(\beta') \cap [Y]$
where $\alpha', \beta' \in K_0(\textit{Vect}(X)) \otimes \mathbf{Q}$
as in Section \ref{section-intersection-regular}.
Setting $c = ch(\alpha')$ and $c' = ch(\beta')$ we find
$\alpha \cdot \beta = c \cap c' \cap [Y]$ by construction.
By Lemma \ref{lemma-lci-gysin-commutes} we know that $f^!$
commutes with both $c$ and $c'$. Hence
\begin{align*}
f^!(\alpha \cdot \beta)
& =
f^!(c \cap c' \cap [Y]) \\
& =
c \cap c' \cap f^![Y] \\
& =
c \cap c' \cap [X] \\
& =
(c \cap [X]) \cdot (c' \cap [X]) \\
& =
(c \cap f^![Y]) \cdot (c' \cap f^![Y]) \\
& =
f^!(\alpha) \cdot f^!(\beta)
\end{align*}
as desired.
\end{proof}

\end{document}
